% Options for packages loaded elsewhere
\PassOptionsToPackage{unicode}{hyperref}
\PassOptionsToPackage{hyphens}{url}
\documentclass[
]{article}
\usepackage{xcolor}
\usepackage{amsmath,amssymb}
\setcounter{secnumdepth}{-\maxdimen} % remove section numbering
\usepackage{iftex}
\ifPDFTeX
  \usepackage[T1]{fontenc}
  \usepackage[utf8]{inputenc}
  \usepackage{textcomp} % provide euro and other symbols
\else % if luatex or xetex
  \usepackage{unicode-math} % this also loads fontspec
  \defaultfontfeatures{Scale=MatchLowercase}
  \defaultfontfeatures[\rmfamily]{Ligatures=TeX,Scale=1}
\fi
\usepackage{lmodern}
\ifPDFTeX\else
  % xetex/luatex font selection
\fi
% Use upquote if available, for straight quotes in verbatim environments
\IfFileExists{upquote.sty}{\usepackage{upquote}}{}
\IfFileExists{microtype.sty}{% use microtype if available
  \usepackage[]{microtype}
  \UseMicrotypeSet[protrusion]{basicmath} % disable protrusion for tt fonts
}{}
\makeatletter
\@ifundefined{KOMAClassName}{% if non-KOMA class
  \IfFileExists{parskip.sty}{%
    \usepackage{parskip}
  }{% else
    \setlength{\parindent}{0pt}
    \setlength{\parskip}{6pt plus 2pt minus 1pt}}
}{% if KOMA class
  \KOMAoptions{parskip=half}}
\makeatother
\usepackage{longtable,booktabs,array}
\newcounter{none} % for unnumbered tables
\usepackage{calc} % for calculating minipage widths
% Correct order of tables after \paragraph or \subparagraph
\usepackage{etoolbox}
\makeatletter
\patchcmd\longtable{\par}{\if@noskipsec\mbox{}\fi\par}{}{}
\makeatother
% Allow footnotes in longtable head/foot
\IfFileExists{footnotehyper.sty}{\usepackage{footnotehyper}}{\usepackage{footnote}}
\makesavenoteenv{longtable}
\setlength{\emergencystretch}{3em} % prevent overfull lines
\providecommand{\tightlist}{%
  \setlength{\itemsep}{0pt}\setlength{\parskip}{0pt}}
\usepackage{graphicx}
\usepackage{bookmark}
\IfFileExists{xurl.sty}{\usepackage{xurl}}{} % add URL line breaks if available
\urlstyle{same}
\hypersetup{
  hidelinks,
  pdfcreator={LaTeX via pandoc}}

\author{}
\date{}

% md2tex unicode
\DeclareUnicodeCharacter{2212}{\ensuremath{-}}
\DeclareUnicodeCharacter{2192}{\ensuremath{\rightarrow}}
\DeclareUnicodeCharacter{2264}{\ensuremath{\le}}
\DeclareUnicodeCharacter{2265}{\ensuremath{\ge}}
\DeclareUnicodeCharacter{2248}{\ensuremath{\approx}}
\DeclareUnicodeCharacter{207B}{\textsuperscript{-}}
\DeclareUnicodeCharacter{03C4}{\ensuremath{\tau}}
\DeclareUnicodeCharacter{03B8}{\ensuremath{\theta}}
\DeclareUnicodeCharacter{03C6}{\ensuremath{\varphi}}
\DeclareUnicodeCharacter{03BE}{\ensuremath{\xi}}
\DeclareUnicodeCharacter{03C1}{\ensuremath{\rho}}
\DeclareUnicodeCharacter{03C9}{\ensuremath{\omega}}
\DeclareUnicodeCharacter{03C8}{\ensuremath{\psi}}
\DeclareUnicodeCharacter{0394}{\ensuremath{\Delta}}
\DeclareUnicodeCharacter{03B1}{\ensuremath{\alpha}}
\DeclareUnicodeCharacter{03BB}{\ensuremath{\lambda}}
\DeclareUnicodeCharacter{03BC}{\ensuremath{\mu}}
\DeclareUnicodeCharacter{03C3}{\ensuremath{\sigma}}
\DeclareUnicodeCharacter{1E61}{\ensuremath{\dot s}}
\DeclareUnicodeCharacter{010B}{\ensuremath{\dot c}}
\DeclareUnicodeCharacter{2500}{-}
\DeclareUnicodeCharacter{2502}{|}
\DeclareUnicodeCharacter{250C}{+}
\DeclareUnicodeCharacter{2510}{+}
\DeclareUnicodeCharacter{2514}{+}
\DeclareUnicodeCharacter{2518}{+}
\DeclareUnicodeCharacter{251C}{+}
\DeclareUnicodeCharacter{2524}{+}
\DeclareUnicodeCharacter{252C}{+}
\DeclareUnicodeCharacter{2534}{+}
\DeclareUnicodeCharacter{253C}{+}
\DeclareUnicodeCharacter{25BA}{>}
\DeclareUnicodeCharacter{25C4}{<}
\DeclareUnicodeCharacter{25BC}{v}
\DeclareUnicodeCharacter{25B2}{^}
% end md2tex unicode
\begin{document}

\section{Tecniche di controllo di RoTino}\label{tecniche-di-controllo-di-rotino}

Documento tecnico sulle due leggi di controllo del workspace \texttt{rotino\_ws}: \textbf{PID con lo ZMP} (\texttt{rotino\_pid}) e \textbf{MPC + TV-LQR + VMC} (\texttt{rotino\_mpc}). Entrambe agiscono sullo stesso URDF, con la stessa attuazione in coppia a 500 Hz e la stessa libreria di modello (\texttt{rotino\_description}). Il documento descrive il modello dinamico su cui poggiano, la struttura di ogni anello, l'origine numerica dei guadagni e delle costanti, e chiude con una sintesi comparativa. Il confronto misurato in Gazebo si rigenera con la suite di benchmark (\texttt{rotino\_benchmark}, Appendice C).

Scritto il 24/09/2026; aggiornato il 01/10/2026: PID riprogettato con lo ZMP (§3) e rimozione della cascata sliding mode dal progetto.

\subsection{Come leggere le affermazioni}\label{come-leggere-le-affermazioni}

Ogni risultato quantitativo porta una di queste etichette:

\begin{longtable}[]{@{}>{\raggedright\arraybackslash}p{(\linewidth - 4\tabcolsep) * \real{0.5000}}>{\raggedright\arraybackslash}p{(\linewidth - 4\tabcolsep) * \real{0.5000}}@{}}
\toprule\noalign{}
Etichetta & Significato \\
\midrule\noalign{}
\endhead
\textbf{[D]} & Derivato analiticamente nel testo, a partire dal modello. \\
\textbf{[C]} & Calcolato eseguendo il codice del workspace (\texttt{WBRModel}, \texttt{lqr\_gain}, \texttt{UpperBodyMPC}, \texttt{AxisKalman}) con i parametri reali dell'URDF. \\
\textbf{[S-R]} & Simulato su modello ridotto: VL-WIP sagittale non lineare (eq. 4-5 di Cui et al.), integrazione RK4 a 2 kHz, controllo a 500 Hz con gli stessi filtri e saturazioni dei nodi. \\
\textbf{[S-G]} & Misurato in Gazebo per questo documento, con \texttt{rotino\_benchmark campaign} e con la registrazione di \texttt{/rotino/wbr\_state}. \\
\textbf{[E]} & Empirico riportato: misura dichiarata nei commenti del codice, in \texttt{docs/Storico\_lavoro.md} o in \texttt{diagnostics/06\_diagnosi.txt}, non ripetuta qui. \\
\bottomrule\noalign{}
\end{longtable}

Gli script che producono i numeri [C], [S-R] e [S-G] sono in \texttt{docs/verifiche/} (Appendice C).

Riferimento bibliografico principale: Z. Cui, Y. Xin, S. Liu, X. Rong, Y. Li, \emph{Modeling and Control of a Wheeled Biped Robot}, Micromachines 2022, 13, 747 (nel seguito "il paper"). I numeri di equazione citati sono quelli del paper.

\begin{center}\rule{0.5\linewidth}{0.5pt}\end{center}

\subsection{0. Convenzioni}\label{0-convenzioni}

\textbf{Terne.} REP 103: \(x\) in avanti, \(y\) a sinistra, \(z\) in alto. Tutti i giunti (anca, ginocchio, ruota) ruotano attorno a \(+y\). La ruota sinistra sta a \(+y\).

\textbf{Coordinate generalizzate del VL-WIP.} \(\phi = [s,\ \theta,\ \varphi]^T\):

\begin{itemize}
\tightlist
\item \(s\): spostamento dell'asse ruote lungo la direzione di marcia [m];
\item \(\theta\): inclinazione del pendolo equivalente, cioè del baricentro del corpo superiore rispetto all'asse, positiva quando il baricentro è \textbf{davanti} all'asse [rad];
\item \(\varphi\): imbardata [rad].
\end{itemize}

\textbf{Segno delle coppie.} Una coppia positiva sul giunto ruota (asse \(+y\)) fa ruotare la ruota in avanti: il centro ruota avanza con velocità \(\omega r\), perché \(\omega\hat y \times r\hat z = \omega r\,\hat x\). La reazione sul corpo è una coppia \(-\tau\) attorno a \(+y\). Una rotazione positiva attorno a \(+y\) porta \(\hat z\) verso \(+\hat x\), cioè inclina il baricentro in avanti. Ne segue che una coppia positiva alle ruote \textbf{accelera l'asse in avanti e fa ruotare il pendolo all'indietro}: nel modello linearizzato \(b_1 > 0\) e \(b_2 < 0\). PID e MPC inviano le coppie in N·m senza inversioni di segno, e pubblicano l'inclinazione con la stessa convenzione (positiva con il baricentro davanti all'asse).

\textbf{Modo comune e differenziale.} Con \(\tau_l, \tau_r\) le coppie delle due ruote: \[\tau_c = \tfrac12(\tau_l+\tau_r), \qquad \tau_d = \tfrac12(\tau_r-\tau_l), \qquad \tau_l = \tau_c-\tau_d,\ \ \tau_r = \tau_c+\tau_d .\] Il modo comune agisce sul piano sagittale (\(s\), \(\theta\)), il differenziale sull'imbardata.

\begin{center}\rule{0.5\linewidth}{0.5pt}\end{center}

\subsection{1. Il modello dinamico}\label{1-il-modello-dinamico}

\subsubsection{1.1 Parametri fisici}\label{11-parametri-fisici}

Tutti i parametri vengono letti dall'URDF (\texttt{src/rotino\_description/urdf/rotino.urdf.xacro}) da \texttt{WBRModel} (\texttt{src/rotino\_description/rotino\_description/model.py:76}). Nessun valore è copiato a mano nei controllori, con l'unica eccezione del PID (§3.1). \textbf{[C]}

\begin{longtable}[]{@{}>{\raggedright\arraybackslash}p{(\linewidth - 6\tabcolsep) * \real{0.3333}}>{\raggedright\arraybackslash}p{(\linewidth - 6\tabcolsep) * \real{0.3333}}>{\raggedright\arraybackslash}p{(\linewidth - 6\tabcolsep) * \real{0.3333}}@{}}
\toprule\noalign{}
Simbolo & Valore & Origine \\
\midrule\noalign{}
\endhead
\(m_w\) & 0,38 kg & ruota 0,32 + mozzo 0,06 \\
\(I_w\) & \(2{,}595\cdot10^{-3}\) kg·m² & \(i_{yy}\) della ruota, include l'armatura MuJoCo di 0,002 \\
\(r\) & 0,06 m & raggio del cilindro di collisione \\
\(d\) & 0,294 m & distanza fra i centri ruota \\
\(m_b\) & 3,53 kg & massa totale (4,29 kg) meno le due ruote \\
\(m_3\) & 2,85 kg & torso: \texttt{base\_link} 2,50 + zavorra 0,35 \\
\(m_2,\ m_1\) & 0,18 kg, 0,16 kg & coscia, stinco (per gamba) \\
\(l_1 = l_2\) & 0,1300 m & \(\sqrt{2}\cdot 0{,}0919\): segmenti inclinati di 45° \\
\(L_{max}\) & 0,2599 m & \(l_1+l_2\) \\
\(\mu\) & 2,2 & \texttt{mu1} delle ruote (DART combina con il minimo; il suolo ha 2,2) \\
\(\tau_{w,max}\) & 18 N·m & limite URDF della ruota (motore MuJoCo, gear 18) \\
\(\tau_{leg,max}\) & 60 N·m & limite URDF di anca e ginocchio \\
\(b_{leg}\) & 0,8 N·m·s/rad & smorzamento viscoso URDF di anca e ginocchio \\
\bottomrule\noalign{}
\end{longtable}

Tre aspetti del robot pesano sul controllo:

\begin{enumerate}
\tightlist
\item \textbf{Rapporto di massa.} Le ruote sono l'18\% della massa totale e, soprattutto, \(2I_w/r^2 = 1{,}44\) kg: l'inerzia rotazionale delle ruote, riportata alla traslazione, vale il 41\% di \(m_b\). Per questo il termine in \(I_w\) nei coefficienti (eq. 14) non è trascurabile.
\item \textbf{Baricentro del torso avanzato.} La zavorra è montata a \((+0{,}045;\ 0;\ -0{,}045)\) m rispetto alla \texttt{base\_link} e il box del torso ha il baricentro a \(x = +0{,}025\) m. Il baricentro del torso sta a \((+0{,}0275;\ -0{,}0055)\) m dall'anca.
\item \textbf{Inerzia del corpo superiore piccola rispetto alla Tabella 1 del paper.} Il paper assume \(I_y = \tfrac13 m_b l^2\) (asta omogenea incernierata all'estremo). Per RoTino l'inerzia reale attorno al baricentro, calcolata dai tensori dell'URDF, è \(I_y = 0{,}0174\) kg·m² nella posa nominale, contro \(\tfrac13 m_b l^2 = 0{,}0311\) kg·m². Il valore della Tabella 1 sarebbe 1,8 volte troppo grande. L'MPC e il feedforward del PID usano il valore reale (§1.5).
\end{enumerate}

\subsubsection{1.2 Cinematica sagittale della gamba}\label{12-cinematica-sagittale-della-gamba}

Ogni gamba è una catena planare di due link con giunti rotoidali attorno a \(y\). Nella posa di riferimento (\(q_{hip} = q_{knee} = 0\)) il ginocchio sta a \((-a, -a)\) dall'anca e la ruota a \((+a, -a)\) dal ginocchio, con \(a = 0{,}0919\) m. La ruota cade quindi esattamente sotto l'anca, \(0{,}1838\) m più in basso.

\texttt{leg\_fk} (\texttt{model.py:211}) calcola la posizione del centro ruota nella terna \texttt{base\_link} e lo Jacobiano rispetto a \((q_{hip}, q_{knee})\). Con \(R_y(q)\) la rotazione attorno a \(y\) proiettata sul piano \((x,z)\): \[p_f(q) = p_{hip} + R_y(q_{hip})\,p_{k} + R_y(q_{hip}+q_{knee})\,p_{w}, \qquad J = \frac{\partial p_f}{\partial q} .\] Nella posa nominale \textbf{[C]}: \[J_0 = \begin{bmatrix} -0{,}1838 & -0{,}0919 \\ 0 & -0{,}0919\end{bmatrix}, \qquad \det J_0 = 0{,}01689\ \text{m}^2, \qquad J_0^{-T} = \begin{bmatrix} -5{,}44 & 0 \\ 5{,}44 & -10{,}88\end{bmatrix}\ \text{m}^{-1}.\] \(J_0^{-T}\) converte coppie di giunto in forze al centro ruota: 1 N·m al ginocchio corrisponde a circa 10,9 N verticali.

**Vincolo \(q_{hip} = -q_{knee}/2\).** Con questo vincolo la ruota resta sulla verticale dell'anca e il torso resta orizzontale. La tabella dell'altezza usata dall'MPC per schedulare i guadagni e per il salto campiona 25 pose con \(q_{hip} \in [-0{,}45;\ 0{,}45]\) e \(q_{knee} = -2q_{hip}\) (\texttt{controller.py:169-175} in entrambi i package).

\subsubsection{1.3 Centroide equivalente (eq. 2-3)}\label{13-centroide-equivalente-eq-2-3}

Il corpo superiore (torso, zavorra, cosce, stinchi) viene ridotto a un punto materiale di massa \(m_b\) nel suo baricentro. Nella terna dell'asse: \[{}^wP_C = \frac{\sum_i m_i\,{}^wP_{Ci}(q)}{\sum_i m_i} = [S_C,\ Z_C]^T, \qquad l = \sqrt{S_C^2+Z_C^2}, \qquad \theta = \operatorname{atan2}(S_C, Z_C).\] \texttt{equivalent\_centroid} (\texttt{model.py:151}) calcola anche l'inerzia \(I_y\) attorno al baricentro, sommando i tensori dei link ruotati in terna mondo e i termini di trasporto di Steiner. \textbf{[C]}

\begin{longtable}[]{@{}>{\raggedright\arraybackslash}p{(\linewidth - 16\tabcolsep) * \real{0.1250}}>{\raggedright\arraybackslash}p{(\linewidth - 16\tabcolsep) * \real{0.1250}}>{\raggedright\arraybackslash}p{(\linewidth - 16\tabcolsep) * \real{0.1250}}>{\raggedright\arraybackslash}p{(\linewidth - 16\tabcolsep) * \real{0.1250}}>{\raggedright\arraybackslash}p{(\linewidth - 16\tabcolsep) * \real{0.1250}}>{\raggedright\arraybackslash}p{(\linewidth - 16\tabcolsep) * \real{0.1250}}>{\raggedright\arraybackslash}p{(\linewidth - 16\tabcolsep) * \real{0.1250}}>{\raggedright\arraybackslash}p{(\linewidth - 16\tabcolsep) * \real{0.1250}}@{}}
\toprule\noalign{}
\(q_{hip}\) & \(q_{knee}\) & \(z_b\) [m] & \(S_C\) [m] & \(Z_C\) [m] & \(l\) [m] & \(\theta\) a torso orizzontale & \(I_y\) [kg·m²] \\
\midrule\noalign{}
\endhead
−0,45 & +0,90 & 0,2454 & +0,0180 & 0,2180 & 0,2188 & +4,73° & 0,0201 \\
−0,30 & +0,60 & 0,2299 & +0,0163 & 0,2040 & 0,2046 & +4,58° & 0,0194 \\
−0,15 & +0,30 & 0,2092 & +0,0147 & 0,1852 & 0,1858 & +4,55° & 0,0184 \\
\textbf{0} & \textbf{0} & \textbf{0,1838} & \textbf{+0,0133} & \textbf{0,1622} & \textbf{0,1627} & \textbf{+4,69°} & \textbf{0,0174} \\
+0,15 & −0,30 & 0,1543 & +0,0121 & 0,1354 & 0,1359 & +5,10° & 0,0163 \\
+0,30 & −0,60 & 0,1213 & +0,0111 & 0,1055 & 0,1061 & +6,00° & 0,0153 \\
+0,45 & −0,90 & 0,0856 & +0,0103 & 0,0731 & 0,0738 & +8,06° & 0,0145 \\
\bottomrule\noalign{}
\end{longtable}

\(z_b\) è la distanza verticale anca-asse. La colonna \(\theta\) è l'inclinazione \textbf{a torso orizzontale}: con il robot in quella posa il baricentro sta 1,0-1,8 cm davanti all'asse. Il robot in equilibrio non può quindi avere il torso orizzontale. In equilibrio \(\theta = 0\), cioè il torso è ruotato all'indietro di circa 4,7° nella posa nominale. Questo spiega il \texttt{PITCH\_OFFSET} del PID (§3.2).

\textbf{Equivalenza fra baricentro totale e baricentro del corpo superiore.} Il PID misura l'inclinazione del baricentro dell'\textbf{intero} robot, l'MPC quella del \textbf{solo corpo superiore}. Le due coincidono. I baricentri delle ruote giacciono sull'asse, quindi rispetto al punto medio dell'asse \(m_{tot}\,\Delta x_{tot} = m_b S_C\) e \(m_{tot}\,\Delta z_{tot} = m_b Z_C\). Il rapporto \(\Delta x/\Delta z\) e quindi l'angolo sono identici. \textbf{[D]}

\subsubsection{1.4 Il VL-WIP non lineare (eq. 4-5)}\label{14-il-vl-wip-non-lineare-eq-4-5}

Il robot viene ridotto a un pendolo inverso su ruote a lunghezza variabile (VL-WIP). Il pendolo ha massa \(m_b\), lunghezza \(l\) e inerzia \(I_y\) attorno al proprio baricentro; lo portano due ruote di massa \(m_w\) e inerzia \(I_w\). Dalle equazioni di Eulero-Lagrange, a \(l\) congelato: \[M(\phi)\ddot\phi + C(\phi,\dot\phi) = B\,\tau_w, \qquad \tau_w = [\tau_l,\ \tau_r]^T\] \[M = \begin{bmatrix} m_b+2m_w+\frac{2I_w}{r^2} & m_b l\cos\theta & 0\\ m_b l\cos\theta & m_b l^2 + I_y & 0 \\ 0 & 0 & \frac{d^2 m_w}{2}+\frac{d^2 I_w}{2r^2}+I_z\end{bmatrix},\quad C = \begin{bmatrix} -m_b l\dot\theta^2\sin\theta \\ -m_b g l \sin\theta \\ 0\end{bmatrix},\quad B = \begin{bmatrix} r^{-1} & r^{-1}\\ -1 & -1 \\ \frac{d}{2r} & -\frac{d}{2r}\end{bmatrix}.\]

Significato dei termini:

\begin{itemize}
\tightlist
\item \(M_{11}\) è la massa traslante equivalente: corpo, due ruote e l'inerzia di rotolamento delle ruote \(2I_w/r^2\).
\item \(M_{12} = m_b l\cos\theta\) è l'accoppiamento inerziale: accelerare l'asse fa ruotare il pendolo e viceversa. È l'origine della fase non minima (§1.6).
\item \(C_2 = -m_b g l\sin\theta\) è la coppia di gravità che rende instabile l'equilibrio.
\item La riga 2 di \(B\) vale \(-1\): la coppia motore agisce sul pendolo con segno opposto a quello con cui fa girare la ruota (reazione).
\item L'imbardata è disaccoppiata: \(M\) è diagonale a blocchi e \(C_3 = 0\).
\end{itemize}

Nel paper la terza riga di \(B\) è \(\left[\frac{d}{2r},\ -\frac{d}{2r}\right]\) e porta a \(B(l)\) con \([b_3,\ -b_3]\). Nel workspace la ruota sinistra è a \(+y\), quindi una coppia in avanti sulla sinistra fa girare il robot in senso orario: \texttt{vlwip\_matrices} usa \([-b_3,\ b_3]\) (\texttt{model.py:207-208}). È un cambio di convenzione, non di fisica.

Il modello trascura: la dinamica di \(l\) (\(\dot l\), \(\ddot l\) e i relativi termini di Coriolis), la rotazione del torso rispetto al pendolo equivalente, la cedevolezza delle gambe, lo slittamento e il distacco delle ruote.

\subsubsection{1.5 Linearizzazione (eq. 13-14)}\label{15-linearizzazione-eq-13-14}

Attorno a \(\theta = 0\), \(\dot\theta = 0\), con \(\sin\theta\approx\theta\), \(\cos\theta\approx 1\) e \(\dot\theta^2\approx 0\), e con \(X = [s,\ \theta,\ \varphi,\ \dot s,\ \dot\theta,\ \dot\varphi]^T\) e \(U = [\tau_l,\ \tau_r]^T\): \[\dot X = A(l)X + B(l)U,\qquad \ddot s = a_1\theta + b_1(\tau_l+\tau_r),\quad \ddot\theta = a_2\theta + b_2(\tau_l+\tau_r),\quad \ddot\varphi = b_3(\tau_r-\tau_l)\] con, posto \(\Delta = 2I_w(I_y+m_bl^2) + \big(2l^2m_bm_w + I_y(m_b+2m_w)\big)r^2\): \[a_1 = -\frac{g\,l^2m_b^2r^2}{\Delta},\quad a_2 = \frac{g\,l\,m_b\big(2I_w+(m_b+2m_w)r^2\big)}{\Delta},\quad b_1 = \frac{r\big(I_y+l\,m_b(l+r)\big)}{\Delta},\quad b_2 = -\frac{2I_w+r\big(l\,m_b+(m_b+2m_w)r\big)}{\Delta}\] \[b_3 = \frac{d\,r}{2I_zr^2+d^2(I_w+m_wr^2)} .\]

\textbf{Verifica.} Le espressioni di \texttt{vlwip\_coefficients} (\texttt{model.py:180-195}) coincidono con l'inversione simbolica di \(M\) linearizzata applicata a \([u/r,\ -u + m_bgl\theta]^T\), con \(u = \tau_l+\tau_r\): la differenza, calcolata con SymPy, è identicamente zero. \(b_1\) e \(b_2\) moltiplicano \textbf{ciascuna} coppia di ruota, cioè la loro somma. \textbf{[D]+[C]}

Coefficienti lungo la griglia di pose, con \(I_y\) reale \textbf{[C]}:

\begin{longtable}[]{@{}>{\raggedright\arraybackslash}p{(\linewidth - 12\tabcolsep) * \real{0.1667}}>{\raggedright\arraybackslash}p{(\linewidth - 12\tabcolsep) * \real{0.1667}}>{\raggedright\arraybackslash}p{(\linewidth - 12\tabcolsep) * \real{0.1667}}>{\raggedright\arraybackslash}p{(\linewidth - 12\tabcolsep) * \real{0.1667}}>{\raggedright\arraybackslash}p{(\linewidth - 12\tabcolsep) * \real{0.1667}}>{\raggedright\arraybackslash}p{(\linewidth - 12\tabcolsep) * \real{0.1667}}@{}}
\toprule\noalign{}
\(l\) [m] & \(a_1\) [m/s²] & \(a_2\) [s⁻²] & \(b_1\) [m/(s²·N·m)] & \(b_2\) [1/(s²·N·m)] & \(\sqrt{a_2}\) [rad/s] \\
\midrule\noalign{}
\endhead
0,2188 & −12,02 & 89,17 & 8,055 & −38,20 & 9,44 \\
0,2046 & −11,73 & 93,07 & 8,040 & −40,72 & 9,65 \\
0,1858 & −11,29 & 98,64 & 8,007 & −44,57 & 9,93 \\
\textbf{0,1627} & \textbf{−10,60} & \textbf{105,82} & \textbf{7,933} & \textbf{−50,15} & \textbf{10,29} \\
0,1359 & −9,54 & 113,94 & 7,762 & −57,98 & 10,67 \\
0,1061 & −7,86 & 120,35 & 7,379 & −68,44 & 10,97 \\
0,0738 & −5,32 & 116,94 & 6,564 & −80,40 & 10,81 \\
\bottomrule\noalign{}
\end{longtable}

\(b_3 = 34{,}88\) rad/(s²·N·m) con \(I_z = 0{,}0227\) kg·m² della posa nominale, tenuto costante.

Con l'\(I_y\) della Tabella 1 (\(\tfrac13 m_b l^2\)), in posa nominale si otterrebbe \(a_2 = 84{,}0\) e \(b_2 = -39{,}8\): 20\% di instabilità in meno e 21\% di autorità in meno rispetto al robot reale. Un controllore che inverte \(b_2\) sbaglierebbe la coppia di quel fattore.

\subsubsection{1.6 Proprietà strutturali}\label{16-proprietà-strutturali}

\textbf{Instabilità.} Il sottosistema \((\theta, \dot\theta)\) ad anello aperto ha i poli \(\pm\sqrt{a_2}\). In posa nominale il polo instabile è a \(+10{,}29\) rad/s: costante di tempo di divergenza 97 ms. Accorciando le gambe il polo sale fino a circa 11 rad/s. \textbf{[C]}

\textbf{Fase non minima.} Dalla coppia totale \(u\) allo spostamento dell'asse, trasformando con Laplace le due equazioni sagittali: \[\frac{S(\sigma)}{U(\sigma)} = \frac{b_1\sigma^2 + (a_1b_2 - a_2b_1)}{\sigma^2(\sigma^2 - a_2)} .\] Lo zero sta in \(\sigma^2 = (a_2b_1 - a_1b_2)/b_1\). In posa nominale \((105{,}82\cdot7{,}933 - 10{,}60\cdot50{,}15)/7{,}933 = 38{,}8\), quindi **zeri reali in \(\pm 6{,}23\) rad/s\textbf{, di cui uno a parte reale positiva (da 5,67 a 7,20 rad/s lungo la griglia). }[D]+[C]**

In termini fisici: per accelerare in avanti a regime il baricentro deve stare davanti all'asse. Per portarcelo le ruote devono prima arretrare. Una coppia positiva produce subito \(\ddot s = b_1u > 0\), ma ruota il pendolo all'indietro (\(b_2 < 0\)); la gravità prende poi il sopravvento e il moto netto si inverte. Nessun controllore causale può seguire un gradino di velocità senza questa sottoelongazione iniziale.

\textbf{Guadagno statico inclinazione → accelerazione.} Se il pendolo è tenuto a inclinazione costante (\(\ddot\theta = 0\)), la coppia necessaria è \(u = -a_2\theta/b_2\) e l'asse accelera con \[\ddot s = \Big(a_1 - \frac{b_1a_2}{b_2}\Big)\theta \equiv g_{st}\,\theta,\qquad g_{st} = 6{,}14\ \mathrm{m/(s^2\cdot rad)}\ \text{in posa nominale}.\] \textbf{[D]+[C]} È il meccanismo su cui si reggono il PID, che comanda l'offset fra baricentro e ZMP (§3.2), e l'MPC (§4.5): \textbf{l'inclinazione è l'ingresso effettivo della traslazione}. Il suo limite superiore fissa la massima accelerazione ottenibile. Per l'inclinazione costante di 10,48°, il limite dell'MPC, \(g_{st}\tan\theta\) vale 1,13 m/s². Il PID limita l'accelerazione comandata a 3 m/s². La formula usa la tangente perché così è definito \(\theta_{ref}\) nell'MPC; per angoli piccoli coincide con \(g_{st}\theta\).

\textbf{Disaccoppiamento.} \(A\) e \(B\) sono a blocchi: sagittale \((s, \theta, \dot s, \dot\theta)\) su \(\tau_l+\tau_r\), imbardata \((\varphi, \dot\varphi)\) su \(\tau_r-\tau_l\). Nel modello lineare il bilanciamento e la sterzata sono indipendenti.

\subsubsection{1.7 Il corpo superiore come massa concentrata (eq. 6-9)}\label{17-il-corpo-superiore-come-massa-concentrata-eq-6-9}

Per pianificare lo spostamento del baricentro il paper usa un secondo modello. Il corpo superiore è una massa \(m_b\) ad altezza \(h\), spostata di \(\Delta s\) in orizzontale rispetto al punto di contatto. La condizione di non ribaltamento richiede che la risultante di gravità e forza d'inerzia passi per il punto di appoggio (eq. 6): \[\frac{F_s}{\Delta s} = \frac{F_z}{h},\qquad F_s = m\ddot s,\quad F_z = m(g+\ddot z)\quad\Longrightarrow\quad \ddot s = \frac{g+\ddot z}{h}\,\Delta s \quad\text{(eq. 7)}\] In verticale la massa è spinta dalla forza delle gambe (eq. 8): \[\ddot z = \frac{F_z}{m_b} - g .\] Lo stato \(x = [s,\ \dot s,\ z,\ \dot z,\ -g]\) con ingresso \(u = [\Delta s,\ F_z]\) dà l'eq. 9, che \texttt{upper\_body\_matrices} (\texttt{model.py:225}) riproduce. Per \(\ddot z\) noto, l'eq. 7 è un pendolo inverso lineare con ingresso la posizione del centro di pressione relativa al baricentro. Il tempo caratteristico è \(\sqrt{h/g} = 0{,}129\) s a \(h = 0{,}162\) m. \textbf{[D]}

I due modelli si completano. Il VL-WIP descrive come le ruote tengono in equilibrio il pendolo; il modello a massa concentrata descrive dove conviene mettere il baricentro rispetto alle ruote, e con quale forza verticale, per muovere il corpo. \(\Delta s\) del secondo modello e \(\theta\) del primo sono legati da \(\theta = \operatorname{atan2}(\Delta s, z)\) (§4.3).

\subsubsection{1.8 Tempo discreto, filtri e ritardi}\label{18-tempo-discreto-filtri-e-ritardi}

\begin{itemize}
\tightlist
\item \textbf{Frequenze.} Fisica a 2 kHz (\texttt{max\_step\_size} 0,5 ms), \texttt{controller\_manager} a 500 Hz, odometria a 500 Hz. I tre nodi eseguono un passo per ogni messaggio di \texttt{/joint\_states}: \(\Delta t = 2\) ms, misurato costante \textbf{[E]}. L'MPC gira ogni cinque passi (100 Hz).
\item \textbf{Filtri del primo ordine} \(y_k = \alpha y_{k-1} + (1-\alpha)x_k\) a 500 Hz. La costante di tempo equivalente è \(\tau = -\Delta t/\ln\alpha\): con \(\alpha = 0{,}8\) (su \(\dot\theta\) e \(\dot s\) nell'MPC) \(\tau = 8{,}96\) ms; con \(\alpha = 0{,}9\) (su \(\Delta s\) nell'MPC) \(\tau = 19{,}0\) ms. \textbf{[D]} Alla frequenza del polo instabile (10,3 rad/s) il filtro su \(\dot\theta\) introduce 5,3° di ritardo di fase. A 70 rad/s ne introduce 32°.
\item **\(\dot\theta\)** è ottenuta per differenza finita di \(\theta\) (calcolato dalla cinematica e dall'IMU) e poi filtrata. Nessun controllore usa direttamente la velocità di beccheggio del giroscopio per \(\dot\theta\): \(\theta\) è l'angolo del baricentro rispetto all'asse, che cambia anche per il moto delle gambe.
\end{itemize}

\begin{center}\rule{0.5\linewidth}{0.5pt}\end{center}

\subsection{2. Stima dello stato (MPC)}\label{2-stima-dello-stato-mpc}

L'MPC usa lo stimatore descritto qui (\texttt{rotino\_mpc/controller.py:504-574}). Il PID \textbf{non} lo usa: legge posa e velocità del torso dalla ground truth del simulatore (\texttt{/rotino/odom}) e ne deriva stati senza filtro. Nel confronto è una differenza di condizioni, non di legge di controllo, e va tenuta presente.

\subsubsection{2.1 Grandezze misurate}\label{21-grandezze-misurate}

\begin{itemize}
\tightlist
\item \textbf{IMU} (500 Hz): quaternione di orientazione \(R\), velocità angolare \(\omega_b\), forza specifica \(f_b\). L'accelerazione in terna mondo è \(a_w = Rf_b - [0,0,g]^T\).
\item \textbf{Encoder} su anche, ginocchia e ruote: posizioni e velocità.
\item \textbf{Sensori di contatto} delle ruote: forza di contatto con il suolo, usata solo per lo stato di appoggio.
\end{itemize}

La direzione di marcia è la proiezione orizzontale normalizzata della prima colonna di \(R\); l'imbardata è \(\varphi = \operatorname{atan2}(h_y, h_x)\).

\subsubsection{2.2 Filtro di Kalman lineare (eq. 20-21)}\label{22-filtro-di-kalman-lineare-eq-20-21}

Per ciascun asse mondo, indipendentemente, lo stato è \([p_b,\ v_b]\) del torso (origine della \texttt{base\_link}). Il modello di processo è un doppio integratore guidato dall'accelerazione IMU: \[\begin{bmatrix}p\\v\end{bmatrix}_{k+1} = \begin{bmatrix}1&\Delta t\\0&1\end{bmatrix}\begin{bmatrix}p\\v\end{bmatrix}_k + \begin{bmatrix}\Delta t^2/2\\ \Delta t\end{bmatrix}a_{w,k} + w_k,\qquad Q = q_a\,\Gamma\Gamma^T,\ \ \Gamma = [\Delta t^2/2,\ \Delta t]^T .\] L'osservazione è la posizione e la velocità del torso ricostruite dall'odometria ruote e dalla cinematica delle gambe: \[y = \begin{bmatrix}P_w + {}^wP_b\\ V_w + {}^wV_b\end{bmatrix},\qquad {}^wP_b = -R\,p_{axle}^{b},\qquad {}^wV_b = -\omega_w\times(R\,p_{axle}^b) - R\,\dot p_{axle}^b .\] \(p_{axle}^b\) è il punto medio dei due centri ruota nella terna \texttt{base\_link}, dato da \texttt{leg\_fk}, e \(\dot p^b_{axle}\) la sua velocità \(J\dot q\).

\textbf{Odometria.} La velocità assoluta di rotazione di ciascuna ruota attorno a \(y\) è la somma delle velocità lungo la catena seriale, perché tutti gli assi sono paralleli: \[\omega_{ruota} = \dot q_{ruota} + \omega_{pitch} + \dot q_{hip} + \dot q_{knee} .\] Con puro rotolamento \(V_w = \tfrac r2(\omega_l+\omega_r)\,\hat h\). \(P_w\) è l'integrale di \(V_w\) con quota fissata a \(r\). Usare solo \(\dot q_{ruota}\) attribuirebbe alla traslazione ogni oscillazione di beccheggio e ogni movimento delle gambe.

\textbf{Equazioni complete del filtro.} L'algoritmo esegue a 500 Hz le classiche equazioni del Filtro di Kalman Discreto Lineare. Avendo osservazioni dirette su tutto lo stato (\(H=I\)), la formulazione si semplifica. La fase di predizione (tramite IMU) aggiorna stato e covarianza: \[x_{k|k-1} = F x_{k-1|k-1} + \Gamma a_{w,k}, \qquad P_{k|k-1} = F P_{k-1|k-1} F^T + Q\] La fase di correzione (tramite Odometria) fonde la predizione con la misura \(y_k\): \[K_k = P_{k|k-1}(P_{k|k-1} + R)^{-1}\] \[x_{k|k} = x_{k|k-1} + K_k (y_k - x_{k|k-1}), \qquad P_{k|k} = (I - K_k)P_{k|k-1}\] L'innovazione \((y_k - x_{k|k-1})\) rappresenta lo scarto tra l'odometria misurata e la propagazione dell'IMU, compensata dinamicamente dal guadagno \(K_k\).

\textbf{Parametri e guadagno a regime.} \(q_a = 0{,}5\) (m/s²)², \(R = \operatorname{diag}(10^{-4}\ \text{m}^2,\ 10^{-3}\ \text{m}^2/\text{s}^2)\). Il guadagno a regime per passo è \textbf{[C]} \[K_\infty = \begin{bmatrix}0{,}0070 & 0{,}0017\\ 0{,}0168 & 0{,}0434\end{bmatrix}.\] Gli autovalori di \((I-K_\infty)F\) sono 0,9929 e 0,9567, cioè costanti di tempo di 280 ms e 45 ms. L'IMU domina sotto i 45 ms e l'odometria sopra i 280 ms. Nel mezzo le due fonti si fondono. Il filtro serve a eliminare il rumore di derivazione e i salti dell'odometria, non a correggere derive lente dell'IMU. Queste sono comunque limitate perché l'orientazione arriva già stimata.

\textbf{Gating sul contatto.} La correzione avviene solo in appoggio. In volo l'odometria non ha significato, \(P_w\) viene riallineato alla stima (\(P_w = P_b - {}^wP_b\)) e il filtro procede in sola predizione. Il paper, invece, gonfia la covarianza di osservazione durante il salto (Sez. 4.3); l'implementazione la esclude del tutto.

\textbf{Differenze dal paper.} Il paper usa un solo filtro con \(P_b, V_b\in\mathbb R^2\). Qui ci sono tre filtri scalari indipendenti, uno per asse mondo. L'equivalenza vale perché \(F\), \(Q\) e \(R\) sono diagonali a blocchi per asse.

\subsubsection{2.3 Grandezze derivate}\label{23-grandezze-derivate}

\begin{itemize}
\tightlist
\item \(s\): integrale della proiezione della velocità dell'asse \(V_{axle} = V_b - {}^wV_b\) sulla direzione di marcia.
\item \(\dot s\): la stessa proiezione, filtrata con \(\alpha = 0{,}8\).
\item \(\theta\), \(l\): dal centroide equivalente (§1.3), usando \(R\) per portare in terna mondo il vettore asse→baricentro.
\item \(\dot\theta\): differenza finita filtrata (§1.8).
\end{itemize}

L'MPC usa in più la velocità del \textbf{baricentro} \(V_{com} = V_b + \omega_w\times(Rc_b)\) come stato del proprio modello orizzontale (\texttt{rotino\_mpc/controller.py:570-572}).

\subsubsection{2.4 Stato di appoggio}\label{24-stato-di-appoggio}

Isteresi sulla forza normale totale (spegnimento sotto 1,5 N, accensione sopra 6 N) con conferma temporale (4 ms per la perdita, 10 ms per il ritorno) e scarto dei messaggi più vecchi di 6 ms. Il sensore di Gazebo pubblica solo mentre c'è contatto, quindi l'assenza di messaggi significa assenza di contatto. L'asimmetria dei tempi di conferma rileva in fretta il decollo e scarta i rimbalzi all'atterraggio.

\textbf{La forza di contatto in Gazebo non è misurata.} In tutte le prove eseguite per questo documento (33.000 campioni) la forza normale totale pubblicata su \texttt{/rotino/debug} vale sempre esattamente 34,629 N \textbf{[S-G]}. È \(2\cdot\tfrac12 m_bg\), il valore di ripiego che \texttt{\_contact\_cb} assegna quando un messaggio di contatto non contiene wrench (MPC riga 283-284). Non è il peso totale (42,1 N), che è quello che un sensore reale riporterebbe. Il sensore di contatto di Gazebo, con questa configurazione, segnala quindi la presenza del contatto ma non la sua intensità. Ne seguono due fatti:

\begin{itemize}
\tightlist
\item le soglie di 1,5 N e 6 N si riducono a "messaggio presente o assente";
\item ogni grandezza calcolata dalla forza normale è di fatto costante in appoggio. Per questo lo ZMP va ricostruito dalla dinamica (\texttt{rotino\_description/zmp.py}, \texttt{docs/Studio\_ZMP.md}).
\item Il sensore pubblica a circa 2 kHz: la sottoscrizione deve avere coda 1, altrimenti i messaggi arrivano già vecchi e il robot risulta in volo (\texttt{docs/PID\_ZMP.md} §7).
\end{itemize}

\begin{center}\rule{0.5\linewidth}{0.5pt}\end{center}

\subsection{3. Controllore PID con lo ZMP}\label{3-controllore-pid-con-lo-zmp}

\texttt{src/rotino\_pid/rotino\_pid/controller.py} (nodo ROS) e \texttt{zmp\_balance.py} (legge di controllo pura, testata offline). Nasce dal porting di un controllore MuJoCo (\texttt{RoTino\_Ctrl\_BalJumpAchieve.py}, bilanciamento con salto verticale). Il 01/10/2026 è stato riprogettato attorno allo \textbf{Zero Moment Point}. Progetto, prove e confronto con la versione precedente sono in \texttt{docs/PID\_ZMP.md}; qui si riassume la struttura.

\subsubsection{3.1 Struttura degli anelli}\label{31-struttura-degli-anelli}

\begin{verbatim}
 percorso pianificato ──► anteprima LIPM ──► c_ref, c_ref ─────────────┐
 (S, trapezio, va e vieni, teleop)                                    ▼
 CoM, contatto ──► errore Capture Point ──► PI (DCM) ──► ZMP desiderato ──► PD sull'offset ZMP ──► τ_c (modo comune)
                   ξ = c + ċ/ω                ρ = 0,4      p_des = c − s_des     + feedforward di modello
 traiettoria planare ──► (C) guida laterale → ψ_cmd ──► (B) PID di rotta + attrito ──► τ_d (modo differenziale)
 gambe ◄── (D) PD cartesiano ruota→anca + peso in feedforward ◄── altezza, salto (supervisore a 7 stati)
                                              τ_l = τ_c − τ_d,   τ_r = τ_c + τ_d
\end{verbatim}

\textbf{Sensori.} Come prima, tutto viene dalla ground truth: posa del torso da \texttt{/rotino/odom}, baricentro e ruote dalla cinematica diretta sull'URDF. \texttt{/joint\_states} e \texttt{/rotino/odom} sono accoppiati per timestamp (\texttt{\_try\_control}). Le derivate sono differenze finite non filtrate. È una differenza di condizioni rispetto all'MPC, che usa IMU e Kalman (§2), e va tenuta presente nel confronto.

\subsubsection{3.2 Anello (A): lo ZMP come ingresso}\label{32-anello-a-lo-zmp-come-ingresso}

Su due ruote lo ZMP longitudinale non può muoversi dentro un appoggio: sta sulla linea dei contatti, e sono le ruote a spostarla. Con il modello \emph{cart-table} del LIPM \[\ddot c = \omega^2 (c - p),\qquad \omega^2 = g/h,\] dove \(c\) è il baricentro, \(p\) lo ZMP (cioè il contatto) e \(h\) l'altezza del baricentro, lo ZMP è la leva con cui si governa il baricentro.

\begin{itemize}
\tightlist
\item \textbf{Anello esterno, PI sul Capture Point.} \(\xi = c + \dot c/\omega\) evolve come \(\dot\xi = \omega(\xi - p)\). Imponendo \(\dot\xi = \dot\xi_{ref} - k_\xi e_\xi - k_i\!\int e_\xi\) si ottiene l'offset desiderato fra baricentro e ZMP:
\end{itemize}

\[s_{des} = c - p_{des} = s_{ff} - \rho\,\big[\dot e_c/\omega + (k_\xi/\omega)\,e_\xi + (k_i/\omega)\!\textstyle\int e_\xi\big],\qquad |s_{des}| \le a_{max}/\omega^2 .\]

\begin{itemize}
\tightlist
\item \textbf{Anello interno, PD sull'offset ZMP:} \(\tau_c = \tau_{ff} + K_s(s - s_{des}) + K_{sd}(\dot s - \dot s_{ff})\).
\item \textbf{Feedforward.} \(s_{ff}\) e \(\tau_{ff}\) sono l'inclinazione e la coppia che mantengono l'accelerazione di riferimento a regime, dal VL-WIP linearizzato (§1.5): 21,8 mm di offset per m/s².
\item \textbf{Anteprima LIPM.} Il riferimento del baricentro è la soluzione limitata di \(\ddot c = \omega^2(c - p_{ref})\): il percorso delle ruote filtrato con il nucleo simmetrico \((\omega/2)e^{-\omega|u|}\). Il robot si inclina \textbf{prima} del gradino di accelerazione, come richiede un sistema a fase non minima. Con comandi da dashboard il futuro non è noto: si usa l'accelerazione corrente del riferimento.
\end{itemize}

\textbf{Guadagni} (per ruota): \(K_s = 60\) N·m/m, \(K_{sd} = 5\) N·m·s/m, \(\rho = 0{,}4\), \(k_\xi = 1{,}5\) s⁻¹, \(k_i = 0{,}5\) s⁻², \(a_{max} = 3\) m/s², saturazione 10 N·m.

**Perché \(\rho < 1\).** La cascata LIPM "pura" presuppone un anello interno infinitamente veloce. A 500 Hz, con 1–2 campioni di ritardo e il modo pendolo-ruote a circa 10 rad/s, il termine \(\dot c/\omega\) a piena autorità diventa retroazione positiva sulla velocità delle ruote. I guadagni sono stati scelti per posizionamento robusto dei poli sul modello discretizzato, su dieci casi: massa ±20 \%, anca ±0,3 rad, ritardo 1 o 2 campioni. Poli nominali: \(-0{,}5\); \(-2{,}3\pm0{,}2j\); \(-9{,}8\); \(-52\) s⁻¹. \textbf{Smorzamento minimo 0,69} su tutti i casi \textbf{[C]}.

\subsubsection{3.3 Anelli (B) e (C): rotta e guida planare}\label{33-anelli-b-e-c-rotta-e-guida-planare}

\textbf{(C) Guida laterale} (solo con \texttt{planar\_enable}). Con \(e_\perp\) l'errore laterale dell'asse: \[\psi_{cmd} = \psi_{ref} - \operatorname{atan}(K_{lat}\,e_\perp)\cdot\min\!\big(1,\ |v_{ref}|/0{,}10\big),\qquad K_{lat} = 2\ \text{rad/m}.\] Il fattore di velocità azzera la correzione da fermo, perché un veicolo differenziale annulla l'errore laterale solo muovendosi.

\textbf{(B) Rotta.} PID sempre attivo: mantiene la rotta anche da fermo e nei moti rettilinei. \[\tau_d = 1{,}8\,(\psi_{cmd}-\psi) + 0{,}45\,(\dot\psi_{ref}-\dot\psi) + 0{,}25\tanh(\dot\psi_{ref}/0{,}01) + 0{,}10\,\dot\psi_{ref} + 0{,}6\!\textstyle\int(\psi_{cmd}-\psi),\] saturato a 2,7 N·m. I termini d'attrito e l'integrale sono quelli identificati in Gazebo per l'MPC (§4.5). Senza di essi il robot girava a scatti per lo stick-slip delle ruote, e lo ZMP laterale seguiva quegli scatti.

\subsubsection{3.4 Anello (D): gambe in spazio cartesiano}\label{34-anello-d-gambe-in-spazio-cartesiano}

\[\tau = J^T\big[K_p(p_{des} - p) - K_d\,J\dot q + F_{ff}\big] + b\,\dot q,\qquad F_{ff} = -\tfrac12 m_bg\,\hat z_b,\]

dove \(p\) è il centro ruota rispetto all'anca in \texttt{base\_link} (\texttt{leg\_fk}) e \(b\dot q\) compensa lo smorzamento dei giunti dell'URDF.

\begin{itemize}
\tightlist
\item \textbf{In appoggio:} \(K_p = \operatorname{diag}(1500, 5000)\) N/m, \(K_d = \operatorname{diag}(40, 80)\) N·s/m.
\item \textbf{Nel salto:} (3000, 4000) / (60, 80).
\item \textbf{In volo:} (800, 800) / (20, 20), senza sostegno.
\end{itemize}

In spazio giunti la rigidezza di appoggio vale circa 51 N·m/rad all'anca e 55 N·m/rad al ginocchio, con circa 1 N·m·s/rad di smorzamento. Un passo da 2 ms non può più scavalcare l'inerzia riflessa della gamba. Il \textbf{ciclo limite bang-bang} del PD di giunto precedente (160–200 N·m/rad, \texttt{diagnostics/06\_diagnosi.txt}) scompare, e con esso la saturazione ad hoc dello smorzamento.

\textbf{Inclinazione in curva (opzionale).} \texttt{zmp\_lateral:=true} accorcia la gamba interna e allunga quella esterna, portando il CoM verso l'interno di \(h a_y/g\) (\(\sin\varphi = a_y/g\)), con la stessa anteprima LIPM. È disattivata di default: con le ruote cilindriche dell'URDF il contatto salta da uno spigolo all'altro quando il camber si inverte (\texttt{PID\_ZMP.md} §6).

\subsubsection{3.5 Supervisore del salto}\label{35-supervisore-del-salto}

Sette stati: BALANCE → PRELOAD → THRUST → FLIGHT → LANDING → RECOVERY → SETTLE → BALANCE. I riferimenti di giunto sono quelli del controllore MuJoCo e vengono convertiti in posizioni cartesiane della ruota con \texttt{leg\_fk}.

\begin{longtable}[]{@{}>{\raggedright\arraybackslash}p{(\linewidth - 6\tabcolsep) * \real{0.3333}}>{\raggedright\arraybackslash}p{(\linewidth - 6\tabcolsep) * \real{0.3333}}>{\raggedright\arraybackslash}p{(\linewidth - 6\tabcolsep) * \real{0.3333}}@{}}
\toprule\noalign{}
Stato & Riferimenti di giunto (anca, ginocchio) & Uscita \\
\midrule\noalign{}
\endhead
PRELOAD & da (0, 0) a (+0,12, −0,22) con smoothstep in 0,80 s: accucciata & dopo 0,80 s \\
THRUST & da PRELOAD a (−0,27, +0,44) con \(1-(1-\tau)^3\) in 0,02 s: estensione brusca & decollo, o 0,16 s senza decollo → LANDING \\
FLIGHT & posa congelata al decollo & contatto confermato \\
LANDING & dalla posa di contatto a (+0,08, −0,20) in 0,80 s: gambe cedevoli & dopo 0,80 s \\
RECOVERY & ritorno a (0, 0) in 2,0 s & dopo 2,0 s \\
SETTLE & (0, 0) & assestamento per 0,30 s o 12 s di attesa \\
\bottomrule\noalign{}
\end{longtable}

\begin{itemize}
\tightlist
\item \textbf{Ruote nelle fasi di salto.} La correzione del Capture Point è scalata: 0,5 in PRELOAD, 0,1 in THRUST, 0,3 in LANDING e 0,5 in RECOVERY. L'integrale è congelato. In volo le ruote sono solo smorzate.
\item \textbf{Decollo.} È confermato solo se valgono insieme quattro condizioni: contatto perso, forza sotto soglia, distacco geometrico della ruota e velocità verticale del baricentro sopra 0,08 m/s.
\item \textbf{Salto a comando.} Il salto parte anche da \texttt{/rotino/cmd\_jump}: il robot si ferma, salta sul posto e al termine riprende il comando da dove è atterrato.
\end{itemize}

\begin{center}\rule{0.5\linewidth}{0.5pt}\end{center}

\subsection{4. MPC + TV-LQR + VMC}\label{4-mpc-tv-lqr-vmc}

\texttt{src/rotino\_mpc/rotino\_mpc/controller.py} e \texttt{solvers.py}. Implementa l'architettura della Fig. 4 del paper.

\begin{verbatim}
 riferimenti (s, ṡ, z, ż sull'orizzonte) ──► MPC corpo superiore, 100 Hz ──► Δs, F_z
                                              │ piano del baricentro x_h(t+Δt)
                                              ▼
 stima (s, θ, φ, ṡ, θ, φ) ──► TV-LQR K(l), 500 Hz ──► τ_l, τ_r (ruote)
                                  X_ref = [s_plan − Δs, atan2(Δs, z_ref), φ_ref, ṡ_plan, 0, φ_ref]
 Δs, z_ref, F_z ──────────────► VMC task-space, 500 Hz ──► τ_hip, τ_knee (gambe)
\end{verbatim}

\subsubsection{4.1 Perché questa architettura: le motivazioni del paper}\label{41-perché-questa-architettura-le-motivazioni-del-paper}

Il paper parte da una constatazione (Sez. 1). I controllori per robot bipedi su ruote pubblicati fino ad allora si dividevano in due gruppi:

\begin{itemize}
\tightlist
\item \textbf{pendolo inverso su ruote puro} (LQR con feedforward, UDE, IDA-PBC): semplici, ma ignorano l'influenza del torso sul moto;
\item \textbf{whole-body control} (Ascento, Xin et al.): completi, ma complessi da modellare e da calcolare.
\end{itemize}

La proposta è intermedia: \textbf{disaccoppiare} il robot in due modelli semplici e dare a ciascuno il controllore più adatto.

\begin{enumerate}
\tightlist
\item \textbf{VL-WIP + TV-LQR per l'equilibrio.} L'equilibrio è un problema veloce, instabile e ben descritto da un modello lineare a parametri variabili (la lunghezza \(l\) cambia con l'altezza). Un LQR ricalcolato in funzione di \(l\) dà la miglior controreazione lineare per un compromesso dichiarato fra errore e sforzo di controllo (eq. 15). Può girare alla frequenza più alta perché è una moltiplicazione matrice-vettore.
\item \textbf{Massa concentrata + MPC per il corpo superiore.} Dove mettere il baricentro rispetto alle ruote e quanta forza verticale chiedere alle gambe è un problema più lento, ma \textbf{vincolato}:
\item il baricentro non può scostarsi dal contatto oltre il cono d'attrito (eq. 18: \(|\Delta s|\le\mu h\), approssimazione a piramide del cono);
\item le gambe hanno uno spazio di lavoro finito (\(|\Delta s|\le\sqrt{L_{max}^2-z_b^2}\));
\item la forza verticale è limitata (\(F_{min}\le F_z\le F_{max}\)): \(F_{min} > 0\) mantiene il contatto, \(F_{max}\) rispetta gli attuatori.
\end{enumerate}

L'MPC tratta i vincoli direttamente nell'ottimizzazione invece di saturare a posteriori. Inoltre \textbf{predice} l'evoluzione dello stato su un orizzonte, quindi reagisce ai riferimenti futuri prima che arrivino.

\begin{enumerate}
\tightlist
\item \textbf{VMC per tradurre in coppie di giunto.} Il VMC (eq. 19) realizza \(\Delta s\) e \(F_z\) come molla-smorzatore virtuale al piede più una forza in avanti. Non serve la dinamica inversa completa delle gambe in appoggio.
\item \textbf{Filtro di Kalman lineare} invece dell'EKF di Bloesch et al., ritenuto troppo costoso (Sez. 4.3).
\end{enumerate}

Il paper sostiene l'uso dell'MPC con la letteratura citata nell'introduzione:

\begin{itemize}
\tightlist
\item permette di integrare facilmente i vincoli e di migliorare la stabilità, prevedendo il tempo di volo e la sottoattuazione (rif. 21-22);
\item un controllo con anteprima compensa l'errore di ZMP dovuto alla differenza fra modello semplice e modello multicorpo (Kajita et al., rif. 23);
\item il controllo ZMP con anteprima migliora la reiezione dei disturbi nei bipedi (Wieber, rif. 24), "cosa cruciale anche per i WBR".
\end{itemize}

La motivazione esplicita della Sez. 4.2 è: \emph{"in order to improve the robustness of control, MPC is used to solve the horizontal displacement of the CoM and the vertical force. It can predict the state variation of the WBR in a longer time frame."}

\subsubsection{4.2 Vantaggi concreti nell'implementazione di RoTino}\label{42-vantaggi-concreti-nellimplementazione-di-rotino}

Quanto segue deriva dal codice, non dal paper.

\textbf{(a) Anticipo sui riferimenti in un sistema a fase non minima.} L'MPC riceve i riferimenti \(s(t_i), \dot s(t_i), z(t_i), \dot z(t_i)\) per tutti i 25 nodi dell'orizzonte (\texttt{\_run\_mpc}, riga 703). Il robot deve inclinarsi \emph{prima} di accelerare (§1.6) e l'inclinazione ha un limite (§4.5). Vedere il riferimento con 0,5 s di anticipo permette di distribuire lo spostamento del baricentro e di iniziarlo prima della variazione del riferimento. L'LQR, da solo, reagisce solo all'errore presente. Con un profilo di velocità noto la differenza è strutturale.

\textbf{(b) Vincoli rispettati nel piano, non solo nel comando.} Il limite su \(\Delta s\) entra nella QP. Il piano che l'MPC restituisce, cioè la traiettoria predetta del baricentro, rispetta quindi il limite su tutto l'orizzonte. L'LQR insegue quel piano (\texttt{s\_plan}, righe 642-647) e non il riferimento grezzo: il riferimento dell'anello veloce è già compatibile con i limiti fisici. Il PID invece satura l'offset ZMP desiderato a posteriori (§3.2) e protegge l'integrale con un anti-windup.

\textbf{(c) Un'unica sintesi per la forza verticale.} \(F_z\) viene pianificata con lo stesso modello e con limiti che dipendono dalla fase: in spinta il tetto sale da 3 a 6 volte il peso. Lo stesso blocco gestisce appoggio, variazione di quota e spinta del salto.

\textbf{(d) Separazione dei tempi.} 100 Hz bastano: il tempo caratteristico del modello orizzontale è \(\sqrt{h/g} = 0{,}13\) s, cioè 13 intervalli dell'MPC. La parte instabile e veloce resta all'LQR a 500 Hz.

\textbf{Costi e limiti:}

\begin{itemize}
\tightlist
\item \textbf{Nessuna azione integrale.} Né l'LQR né l'MPC hanno stati integrali o un modello del disturbo. Sul modello ridotto, un disturbo persistente lascia al solo TV-LQR un errore di posizione a regime: 0,65 m con 3 N costanti sul torso, 0,08 m con 0,3 N·m di bias per ruota \textbf{[S-R]}. Per l'MPC completo non l'ho misurato, ma nulla nella sua formulazione annulla quell'errore. Il PID ha invece un integrale sul Capture Point (§3.2).
\item \textbf{Modello di predizione povero.} Il corpo superiore è un punto, senza dinamica di beccheggio. L'MPC non sa che per spostare il baricentro di \(\Delta s\) il pendolo deve prima essere portato all'angolo corrispondente: lo delega all'LQR.
\item \textbf{Costo di calcolo.} Due QP a 25 variabili ogni 10 ms, risolte con 60 iterazioni di gradiente accelerato.
\end{itemize}

\subsubsection{4.3 Differenze fra paper e implementazione}\label{43-differenze-fra-paper-e-implementazione}

\begin{longtable}[]{@{}>{\raggedright\arraybackslash}p{(\linewidth - 6\tabcolsep) * \real{0.3333}}>{\raggedright\arraybackslash}p{(\linewidth - 6\tabcolsep) * \real{0.3333}}>{\raggedright\arraybackslash}p{(\linewidth - 6\tabcolsep) * \real{0.3333}}@{}}
\toprule\noalign{}
Aspetto & Paper & \texttt{rotino\_mpc} \\
\midrule\noalign{}
\endhead
Frequenze (Fig. 4) & MPC 100 Hz, stimatore 500 Hz, TV-LQR e task-space 1000 Hz & MPC 100 Hz, tutto il resto 500 Hz \\
\(I_y\) & \(\tfrac13 m_b l^2\) (Tabella 1) & reale dall'URDF, schedulato con la posa \\
Discretizzazione MPC & Eulero: \(A_k = I + A_c\Delta T\), \(B_k = B_c\Delta T\) & esatta (ZOH) per il doppio integratore: \(B_k = [\Delta T^2/2,\ \Delta T]^T\) \\
Struttura della QP & un problema sullo stato a 5 componenti dell'eq. 9 & due QP disaccoppiate (orizzontale, verticale) \\
Vincolo su \(\Delta s\) & \(\min(\mu h,\ \sqrt{L_{max}^2-z_b^2})\) & come il paper \textbf{più} un tetto fisso di 0,03 m \\
Solutore & QP condensata (rif. 28) & QP condensata, solo vincoli di box, FISTA \\
Feedforward VMC & dinamica inversa in volo & \(-F_z/2\) lungo la verticale e compensazione dello smorzamento URDF; nessuna dinamica inversa \\
Kalman & un filtro 2D & tre filtri scalari (equivalente) \\
Imbardata & solo tramite la riga di \(\varphi\) del TV-LQR & TV-LQR con peso differenziale separato, feedforward d'attrito e integrale \\
\bottomrule\noalign{}
\end{longtable}

Nell'eq. 9 l'unico accoppiamento fra orizzontale e verticale è il coefficiente \((g+\ddot z)/h\). Congelandolo sull'orizzonte con il valore del riferimento, i due blocchi diventano indipendenti e ciascuno è lineare tempo-invariante. Il paper lo tratta implicitamente allo stesso modo, perché anche lì la matrice \(B_c\) contiene \(\ddot z\) e \(h\).

\subsubsection{4.4 TV-LQR (eq. 15-16)}\label{44-tv-lqr-eq-15-16}

\textbf{Formulazione.} \(J = \int_0^\infty(\tilde X^TQ\tilde X + U^TRU)\,dt\), \(U = -K(l)\tilde X\). \(K\) si ottiene dall'equazione algebrica di Riccati, risolta come sottospazio stabile dell'Hamiltoniana (\texttt{care}, \texttt{solvers.py:15}).

\textbf{Pesi} (\texttt{controller.py:41-48}): \[Q = \operatorname{diag}(30,\ 400,\ 80,\ 15,\ 6,\ 2)\quad\text{su}\ (s,\ \theta,\ \varphi,\ \dot s,\ \dot\theta,\ \dot\varphi),\qquad R = T^T\operatorname{diag}(2,\ 100)\,T,\ \ T = \begin{bmatrix}\tfrac12&\tfrac12\\-\tfrac12&\tfrac12\end{bmatrix}.\]

\begin{itemize}
\tightlist
\item **Lettura di \(Q\) alla Bryson.** \(1/\sqrt{q_{ii}}\) è lo scostamento "equivalente" a un'unità di costo: 0,18 m in \(s\), 0,050 rad (2,9°) in \(\theta\), 0,11 rad in \(\varphi\), 0,26 m/s in \(\dot s\), 0,41 rad/s in \(\dot\theta\), 0,71 rad/s in \(\dot\varphi\). L'inclinazione ha la priorità, la posizione è accettata a decine di centimetri: il robot "cede" in traslazione per non cadere.
\item **Lettura di \(R\).** \(T\) trasforma \((\tau_l, \tau_r)\) in \((\tau_c, \tau_d)\), per cui \(U^TRU = 2\tau_c^2 + 100\tau_d^2\). Con \(R = I\) il costo sarebbe \(\tau_l^2+\tau_r^2 = 2\tau_c^2 + 2\tau_d^2\). Il peso sul modo comune coincide quindi con \(R = I\), come dichiara il commento a riga 45, mentre il modo differenziale è penalizzato 50 volte di più. Il modello è disaccoppiato (§1.6) e con questa \(R\) anche il costo lo è: la Riccati si separa in un problema sagittale e uno di imbardata, e i due pesi agiscono ciascuno solo sulla propria parte.
\end{itemize}

\textbf{Perché il differenziale è penalizzato così.} Il commento (righe 42-44) motiva la scelta con una risonanza torsionale delle gambe in appoggio, a circa 76 rad/s (12 Hz): l'inerzia di imbardata del torso sulle molle virtuali del VMC \textbf{[E]}. Un anello di imbardata veloce la eccita in un ciclo limite permanente. Con \(R_d = 100\) l'anello di imbardata ha pulsazione di attraversamento \textbf{15,4 rad/s [C]}, un quinto della risonanza. Il modello VL-WIP non contiene questa risonanza: il vincolo è stato scoperto in simulazione e imposto attraverso il peso.

\textbf{Guadagni e poli} in posa nominale \textbf{[C]}: \[K = \begin{bmatrix} -3{,}873 & -20{,}28 & -0{,}894 & -5{,}637 & -2{,}784 & -0{,}214\\ -3{,}873 & -20{,}28 & +0{,}894 & -5{,}637 & -2{,}784 & +0{,}214\end{bmatrix}\] Poli: \(-179{,}5\); \(-8{,}13\); \(-1{,}08\pm0{,}68j\) (sagittali); \(-7{,}45\pm2{,}62j\) (imbardata). Il polo a \(-179\) rad/s sta oltre la banda del filtro su \(\dot\theta\) (\(1/8{,}96\ \text{ms} = 112\) rad/s). Nel sistema reale quel modo è determinato dal filtro e dal campionamento più che dall'LQR. \(-8{,}1\) è la dinamica di inclinazione. La coppia lenta \(-1{,}08\pm0{,}68j\) (\(\zeta = 0{,}85\)) è il ritorno in posizione: è la dinamica che si osserva dopo una spinta.

**Scheduling su \(l\).** \texttt{LQRSchedule} (\texttt{solvers.py:33}) calcola \(K\) sui 25 punti della griglia di pose, ciascuno con il proprio \(I_y\), e interpola linearmente in \(l\) a ogni passo. Lungo la griglia i poli sagittali restano quasi fermi \textbf{[C]}:

\begin{longtable}[]{@{}>{\raggedright\arraybackslash}p{(\linewidth - 10\tabcolsep) * \real{0.2000}}>{\raggedright\arraybackslash}p{(\linewidth - 10\tabcolsep) * \real{0.2000}}>{\raggedright\arraybackslash}p{(\linewidth - 10\tabcolsep) * \real{0.2000}}>{\raggedright\arraybackslash}p{(\linewidth - 10\tabcolsep) * \real{0.2000}}>{\raggedright\arraybackslash}p{(\linewidth - 10\tabcolsep) * \real{0.2000}}@{}}
\toprule\noalign{}
\(l\) [m] & \(K_\theta\) [N·m/rad] & \(K_{\dot\theta}\) & \(K_{\dot s}\) & poli sagittali \\
\midrule\noalign{}
\endhead
0,219 & 21,22 & 3,13 & 5,51 & \(-139{,}9\); \(-8{,}06\); \(-1{,}15\pm0{,}68j\) \\
0,163 & 20,28 & 2,78 & 5,64 & \(-179{,}5\); \(-8{,}13\); \(-1{,}08\pm0{,}68j\) \\
0,074 & 18,43 & 2,33 & 6,32 & \(-281{,}1\); \(-8{,}16\); \(-0{,}86\pm0{,}63j\) \\
\bottomrule\noalign{}
\end{longtable}

È lo scopo del TV-LQR: comportamento ad anello chiuso invariante rispetto all'altezza. \(K_s = -3{,}873\) non cambia con \(l\) e coincide numericamente con \(\sqrt{q_s/r_c} = \sqrt{30/2}\). Allo stesso modo il guadagno di imbardata vale \(0{,}894 = \sqrt{q_\varphi/r_d} = \sqrt{80/100}\). I poli d'imbardata non dipendono da \(l\) perché \(I_z\) è tenuto costante. L'interpolazione a \(l\) congelato è giustificata finché \(l\) varia lentamente rispetto alla dinamica di anello chiuso. Durante la spinta del salto questo non è garantito, e il modello ignora i termini in \(\dot l\).

\textbf{Saturazione e riferimento.} Coppia per ruota ±10 N·m (contro i 18 dell'URDF), con priorità al modo comune: il differenziale è limitato a ±2,5 N·m e il comune a \(\pm(10-|\tau_d|)\) (righe 658-660). Il riferimento è \[X_{ref} = \big[s_{plan}-\Delta s,\ \operatorname{atan2}(\Delta s, z_{ref}),\ \varphi_{ref},\ \dot s_{plan},\ 0,\ \dot\varphi_{ref}\big].\] \(s_{plan}, \dot s_{plan}\) sono lo stato del baricentro predetto dall'MPC al nodo successivo, interpolato fra due soluzioni (il "virtual CoM" della Fig. 4). L'asse deve stare \(\Delta s\) dietro il baricentro pianificato e il pendolo all'angolo che realizza quello spostamento.

\subsubsection{4.5 Imbardata: feedforward d'attrito e integrale}\label{45-imbardata-feedforward-dattrito-e-integrale}

Oltre alla riga di \(\varphi\) del TV-LQR (righe 650-660): \[\tau_d \mathrel{+}= \underbrace{0{,}25\tanh\!\big(\dot\varphi_{ref}/0{,}01\big) + 0{,}10\,\dot\varphi_{ref}}_{\text{attrito di strisciamento}} + \underbrace{\operatorname{sat}_{0,6}\Big(-0{,}6\!\int\! e_\varphi\,dt\Big)}_{\text{integrale}} .\]

\begin{itemize}
\tightlist
\item Il \textbf{feedforward} compensa la coppia di strisciamento delle ruote in curva, modellata come Coulomb (0,25 N·m) più viscoso (0,10 N·m per rad/s), con valori identificati in Gazebo \textbf{[E]}. La \(\tanh\) con 0,01 rad/s rende continua la funzione segno. Dipende dal \textbf{riferimento} e non dalla misura, quindi non può generare chattering.
\item L'\textbf{integrale} annulla l'errore statico di rotta dovuto all'attrito residuo. Si attiva solo con errore sopra 0,035 rad o durante una sterzata comandata, per evitare la caccia da stick-slip su errori piccoli, ed è saturato a 0,6 N·m.
\item L'errore di rotta è \textbf{saturato a 0,35 rad} e il riferimento teleoperato non può anticipare il robot di più di 0,35 rad (anti-windup del riferimento, righe 334-337).
\end{itemize}

\subsubsection{4.6 MPC del corpo superiore (eq. 17-18)}\label{46-mpc-del-corpo-superiore-eq-17-18}

\textbf{Modelli di predizione} (\texttt{UpperBodyMPC}, \texttt{solvers.py:79}), \(\Delta T = 0{,}02\) s, \(N = 25\) (orizzonte di 0,5 s): \[\text{orizzontale:}\ \begin{bmatrix}s\\\dot s\end{bmatrix}_{k+1} = \begin{bmatrix}1&\Delta T\\0&1\end{bmatrix}\begin{bmatrix}s\\\dot s\end{bmatrix}_k + \begin{bmatrix}\Delta T^2/2\\\Delta T\end{bmatrix}\frac{g+\ddot z_{ref}}{h}\,\Delta s_k\] \[\text{verticale:}\ \begin{bmatrix}z\\\dot z\end{bmatrix}_{k+1} = \begin{bmatrix}1&\Delta T\\0&1\end{bmatrix}\begin{bmatrix}z\\\dot z\end{bmatrix}_k + \begin{bmatrix}\Delta T^2/2\\\Delta T\end{bmatrix}\frac{F_{z,k}}{m_b} - \begin{bmatrix}\Delta T^2/2\\\Delta T\end{bmatrix}g\] Lo stato orizzontale è quello del \textbf{baricentro} (\(s_{com} = s + S_C\), \(\dot s_{com}\) da \(V_{com}\)), il verticale è \(z = Z_C\) con \(\dot z\) da \(V_{com}\).

\textbf{Condensazione.} Con \(x_i = \Phi_i x_0 + \Gamma_i U + C_i\) (\texttt{\_condense}, \texttt{solvers.py:61}) il problema diventa \[\min_U\ \tfrac12 U^THU + g^TU,\quad H = \Gamma^T\bar S\Gamma + W I,\quad g = \Gamma^T\bar S(\Phi x_0 + C - x_{ref}),\qquad lo\le U\le hi .\] Solo vincoli di box, per cui la proiezione è una saturazione componente per componente e basta un gradiente proiettato accelerato (FISTA, \texttt{box\_qp}, \texttt{solvers.py:48}): 60 iterazioni, passo \(1/\lambda_{max}(H)\), convergenza \(O(1/k^2)\). La soluzione precedente, traslata di un nodo, fa da punto di partenza.

\textbf{Pesi.}

\begin{itemize}
\tightlist
\item \emph{Orizzontale}: \(S_h = \operatorname{diag}(50, 20)\) su \((s, \dot s)\), \(W_h = 200\) su \(\Delta s\). Uno spostamento di 1 cm costa \(200\cdot10^{-4} = 0{,}02\) per nodo, quanto un errore di posizione di 2 cm (\(50\cdot 4\cdot10^{-4}\)). L'MPC accetta di sbilanciare il baricentro solo per errori di posizione più grandi di circa il doppio. Con un gradino di 0,1 m sul riferimento il primo \(\Delta s\) vale 0,025 m, sotto il tetto \textbf{[C]}.
\item \emph{Verticale}: \(S_v = \operatorname{diag}(5000, 150)\) su \((z, \dot z)\), \(W_v = 2\cdot10^{-3}\) su \(F_z\). Il peso agisce su \(F_z\) \textbf{assoluto} e non sullo scostamento dal peso, quindi in linea di principio spinge \(F_z\) sotto \(m_bg\). In anello chiuso sul modello perfetto l'effetto è trascurabile: errore di quota a regime 0,01 mm, \(F_z = 34{,}629\) N contro un peso di 34,629 N \textbf{[C]}. \(S_v\) è 100 volte \(S_h\): la quota va tenuta rigidamente perché in appoggio il VMC non ha una molla verticale (§4.7).
\end{itemize}

\textbf{Vincoli.} \[|\Delta s| \le \min\big(\mu z,\ \sqrt{L_{max}^2-z_b^2},\ 0{,}03\big),\qquad 0{,}3\,m_bg\le F_z\le \{3;\ 6\}\,m_bg .\] In posa nominale \(\mu z = 0{,}357\) m e \(\sqrt{L_{max}^2 - z_b^2} = 0{,}184\) m. \textbf{Il vincolo attivo è sempre il tetto di 0,03 m} (\texttt{DS\_MAX\_CAP}), che non viene dal paper. Il limite d'attrito del paper (\(\mu h\) con \(\mu = 2{,}2\)) è irrealistico come limite di inclinazione, perché permetterebbe \(\operatorname{atan}(2{,}2) = 65°\). Il tetto di 3 cm corrisponde a \(\theta_{ref,max} = \operatorname{atan}(0{,}03/0{,}162) = 10{,}48°\) \textbf{[C]}, entro la zona in cui il VL-WIP linearizzato resta credibile. Per il §1.6 questo tetto fissa anche la massima decelerazione a regime: \(g_{st}\cdot0{,}03/0{,}162 = 1{,}13\) m/s². \(F_{min} = 0{,}3\,m_bg\) impedisce all'MPC di pianificare uno scarico che farebbe perdere il contatto.

**Da \(\Delta s\) all'LQR e alle gambe.** \(\Delta s\) viene filtrato (\(\alpha = 0{,}9\), 19 ms) e usato due volte:

\begin{itemize}
\tightlist
\item nell'LQR, come angolo di riferimento \(\theta_{ref} = \operatorname{atan2}(\Delta s, z_{ref})\) e come arretramento dell'asse rispetto al baricentro pianificato;
\item nel VMC, come posizione orizzontale desiderata della ruota rispetto al baricentro, \(p_{d,x} = c_{b,x} - \Delta s\).
\end{itemize}

Le ruote (inclinando il pendolo) e le gambe (spostando il piede sotto il corpo) realizzano lo stesso spostamento in modo coerente.

\subsubsection{4.7 VMC (eq. 19)}\label{47-vmc-eq-19}

Per ciascuna gamba (\texttt{\_vmc}, riga 452): \[\tau = J^T\big[K_p(p_d - p_f) + K_d(v_d - v_f) + F_{ff}\big] + b_{leg}\,\dot q,\qquad F_{ff} = -\tfrac12F_z\,\hat z_b\] con \(p_d = (c_{b,x} - \Delta s,\ c_{b,z} - z_{ref})\) e \(v_d = (0, -\dot z_{ref})\) nella terna \texttt{base\_link}. \(\hat z_b\) è la verticale mondo espressa in terna corpo. \(F_{ff}\) è la forza che ciascun piede esercita sul suolo: metà del sostegno pianificato, diretta verso il basso. Il termine \(b_{leg}\dot q\) annulla lo smorzamento viscoso dei giunti dell'URDF (0,8 N·m·s/rad), che altrimenti si sommerebbe al \(K_d\) virtuale in modo dipendente dalla configurazione.

\textbf{Guadagni per fase} (righe 78-84):

\begin{longtable}[]{@{}>{\raggedright\arraybackslash}p{(\linewidth - 8\tabcolsep) * \real{0.2500}}>{\raggedright\arraybackslash}p{(\linewidth - 8\tabcolsep) * \real{0.2500}}>{\raggedright\arraybackslash}p{(\linewidth - 8\tabcolsep) * \real{0.2500}}>{\raggedright\arraybackslash}p{(\linewidth - 8\tabcolsep) * \real{0.2500}}@{}}
\toprule\noalign{}
Fase & \(K_p\) [N/m] (x, z) & \(K_d\) [N·s/m] (x, z) & Sostegno verticale \\
\midrule\noalign{}
\endhead
HOLD (ancorato) & 1500, 2000 & 40, 50 & molla \\
Appoggio con MPC & \textbf{1500, 0} & 40, 15 & **solo \(F_z\) dall'MPC** \\
Appoggio senza MPC & 1500, 2000 & 40, 50 & molla \\
Volo & 800, 800 & 20, 20 & nessuno (gambe raccolte) \\
\bottomrule\noalign{}
\end{longtable}

La scelta chiave è \(K_{p,z} = 0\) in appoggio. La quota è regolata dall'MPC attraverso \(F_z\) a 100 Hz. Una molla verticale in parallelo sarebbe un secondo regolatore sulla stessa grandezza e si opporrebbe alle variazioni di quota pianificate. Resta uno smorzamento di 15 N·s/m contro le oscillazioni fra due aggiornamenti dell'MPC. In orizzontale la rigidezza di 1500 N/m con 40 N·s/m ha costante di superficie \(K_p/K_d = 37{,}5\) s⁻¹, cioè 27 ms: il piede converge alla posizione desiderata molto più in fretta del moto del baricentro. Come ordine di grandezza, con metà di \(m_b\) per gamba, \(\omega_n = \sqrt{1500/1{,}765} = 29\) rad/s e \(\zeta = 0{,}39\) \textbf{[C]}.

\subsubsection{4.8 Salto}\label{48-salto}

Macchina a stati: BALANCE → PRELOAD → THRUST → FLIGHT → LANDING → BALANCE. Il salto parte solo a robot fermo (\(|\dot\theta| < 0{,}3\) rad/s, \(|\dot s| < 0{,}08\) m/s). Le quote vengono dalla tabella del §1.2 \textbf{[C]}: accucciata con distanza anca-asse 0,12 m (\(Z_C = 0{,}104\) m), obiettivo di spinta 0,24 m (\(Z_C = 0{,}213\) m), volo 0,17 m (\(Z_C = 0{,}150\) m).

\begin{itemize}
\tightlist
\item \textbf{PRELOAD}: smoothstep da \(z_{stand}\) a \(z_{low}\) in 0,8 s.
\item \textbf{THRUST}: accelerazione costante \(a = v_{to}^2/\big(2(z_{high}-z_{low})\big)\), così il baricentro arriva all'estensione con la velocità di decollo richiesta. Con \(v_{to} = 1{,}2\) m/s: \(a = 6{,}62\) m/s² per 0,181 s. \(F_z = m_b(g+a) = 58{,}0\) N \(= 1{,}67\,m_bg\), dentro il tetto di spinta di \(6\,m_bg\) \textbf{[C]}. Il termine \(\ddot z_{ref}\) entra anche nel coefficiente \((g+\ddot z)/h\) del modello orizzontale: durante la spinta lo stesso \(\Delta s\) produce un'accelerazione orizzontale 1,67 volte più grande, e l'MPC ne tiene conto.
\item \textbf{FLIGHT}: ruote frenate (\(-0{,}05\,\dot q\)), gambe portate sotto il baricentro con il VMC di volo, MPC e LQR sospesi.
\item \textbf{LANDING}: all'atterraggio l'MPC viene azzerato, il riferimento di posizione riallineato alla posizione attuale e la quota riportata con smoothstep in 1,0 s, più 1,5 s di assestamento.
\end{itemize}

\begin{center}\rule{0.5\linewidth}{0.5pt}\end{center}

\subsection{5. Sintesi comparativa}\label{5-sintesi-comparativa}

\begin{longtable}[]{@{}>{\raggedright\arraybackslash}p{(\linewidth - 6\tabcolsep) * \real{0.3333}}>{\raggedright\arraybackslash}p{(\linewidth - 6\tabcolsep) * \real{0.3333}}>{\raggedright\arraybackslash}p{(\linewidth - 6\tabcolsep) * \real{0.3333}}@{}}
\toprule\noalign{}
 & PID con ZMP & MPC + TV-LQR + VMC \\
\midrule\noalign{}
\endhead
Stato per il bilanciamento & ground truth, derivate non filtrate & IMU + Kalman + cinematica \\
Legge sulle ruote & PI sul Capture Point → ZMP desiderato → PD sull'offset ZMP & LQR schedulato su \(l\) \\
Uso del modello & feedforward VL-WIP, anteprima LIPM, guadagni per poli robusti & LQR (VL-WIP) + MPC (massa concentrata) \\
Traslazione & anteprima LIPM del percorso + integrale sul Capture Point & MPC con anteprima + LQR sul piano \\
Limite di inclinazione & offset ZMP ≤ \(a_{max}/\omega^2\) (3 m/s²) & \(\Delta s\le\) 3 cm → 10,48° \\
Imbardata & PID + attrito (Coulomb e viscoso) & LQR (\(R_d = 100\)) + attrito + integrale \\
Gambe & PD cartesiano, 1500 / 5000 N/m, peso in feedforward & VMC, \(K_{p,z} = 0\), sostegno da \(F_z\) \\
Azione integrale & Capture Point, imbardata & solo sull'imbardata \\
Poli sagittali (posa nominale) & \(-0{,}5\); \(-2{,}3\pm0{,}2j\); \(-9{,}8\); \(-52\) & \(-8{,}13\); \(-1{,}08\pm0{,}68j\) \\
Comandi da dashboard & sì (1,5 m/s) & sì (0,6 m/s) \\
\bottomrule\noalign{}
\end{longtable}

Il confronto misurato, scenario per scenario, si ottiene con \texttt{ros2 run rotino\_benchmark suite} (Appendice C): per ogni scenario una tabella PID | MPC con la legge migliore su ogni metrica e i grafici sovrapposti.

\begin{center}\rule{0.5\linewidth}{0.5pt}\end{center}

\subsection{Appendice A. Linearizzazione del VL-WIP passo per passo}\label{appendice-a-linearizzazione-del-vl-wip-passo-per-passo}

Dalle eq. 4-5, riga sagittale, con \(u = \tau_l+\tau_r\) e approssimazioni al primo ordine: \[\underbrace{\begin{bmatrix} m_b+2m_w+2I_w/r^2 & m_bl\\ m_bl & m_bl^2+I_y\end{bmatrix}}_{M_0}\begin{bmatrix}\ddot s\\ \ddot\theta\end{bmatrix} = \begin{bmatrix}u/r\\ -u + m_bgl\,\theta\end{bmatrix}.\] Posto \(M_0 = \begin{bmatrix}\alpha&\beta\\\beta&\gamma\end{bmatrix}\), \(\det M_0 = \alpha\gamma - \beta^2\), si ha \[\ddot s = \frac{\gamma\,u/r - \beta(-u+m_bgl\theta)}{\det M_0},\qquad \ddot\theta = \frac{-\beta\,u/r + \alpha(-u+m_bgl\theta)}{\det M_0}.\] Moltiplicando numeratore e denominatore per \(r^2\): \(r^2\det M_0 = \Delta\) del §1.5, e \[a_1 = -\frac{\beta\,m_bgl\,r^2}{\Delta} = -\frac{g\,l^2m_b^2r^2}{\Delta},\qquad a_2 = \frac{\alpha\,m_bgl\,r^2}{\Delta} = \frac{glm_b\big(2I_w + (m_b+2m_w)r^2\big)}{\Delta}\] \[b_1 = \frac{(\gamma/r+\beta)\,r^2}{\Delta} = \frac{r\big(I_y + lm_b(l+r)\big)}{\Delta},\qquad b_2 = -\frac{(\beta/r+\alpha)\,r^2}{\Delta} = -\frac{2I_w + r\big(lm_b+(m_b+2m_w)r\big)}{\Delta}.\] Sono le espressioni dell'eq. 14 e di \texttt{vlwip\_coefficients}. Imbardata: \(M_{33}\ddot\varphi = \frac{d}{2r}(\tau_l-\tau_r)\), da cui \(b_3 = \frac{d/(2r)}{M_{33}} = \frac{dr}{2I_zr^2 + d^2(I_w+m_wr^2)}\).

\subsection{Appendice B. Costanti principali e loro origine}\label{appendice-b-costanti-principali-e-loro-origine}

\begin{longtable}[]{@{}>{\raggedright\arraybackslash}p{(\linewidth - 8\tabcolsep) * \real{0.2500}}>{\raggedright\arraybackslash}p{(\linewidth - 8\tabcolsep) * \real{0.2500}}>{\raggedright\arraybackslash}p{(\linewidth - 8\tabcolsep) * \real{0.2500}}>{\raggedright\arraybackslash}p{(\linewidth - 8\tabcolsep) * \real{0.2500}}@{}}
\toprule\noalign{}
Costante & Valore & File:riga & Origine \\
\midrule\noalign{}
\endhead
\texttt{LQR\_Q} & diag(30, 400, 80, 15, 6, 2) & \texttt{rotino\_mpc/controller.py:41} & taratura; lettura alla Bryson §4.4 \\
\texttt{LQR\_R\_COMMON} / \texttt{\_DIFF} & 2 / 100 & \texttt{rotino\_mpc/controller.py:45-46} & \(R = I\) sul comune; differenziale sotto la risonanza \textbf{[E]} \\
\texttt{WHEEL\_TORQUE\_MAX} & 10 N·m & MPC :49, PID :56 & margine sui 18 N·m dell'URDF \\
\texttt{DIFF\_TORQUE\_MAX} & 2,5 N·m & MPC :50 & priorità al bilanciamento \\
\texttt{YAW\_FRICTION} / \texttt{VISCOUS} & 0,25 N·m / 0,10 N·m·s & MPC :53-54, PID \texttt{YAW\_FRICTION} & identificati in Gazebo \textbf{[E]} \\
\texttt{YAW\_KI}, \texttt{YAW\_I\_MAX}, \texttt{YAW\_I\_DEADBAND} & 0,6; 0,6; 0,035 & MPC :56-58, PID \texttt{YAW\_KI} & integrale lento, anti stick-slip \\
\texttt{THETA\_DOT\_FILTER}, \texttt{S\_DOT\_FILTER} & 0,8 (τ = 9 ms) & MPC :60-61 & ritardo di fase 5,3° a 10 rad/s \\
\texttt{MPC\_HORIZON}, \texttt{MPC\_DT} & 25, 0,02 s & \texttt{rotino\_mpc/controller.py:66-67} & orizzonte 0,5 s ≈ 4 \(\sqrt{h/g}\) \\
\texttt{MPC\_S\_H}, \texttt{MPC\_W\_H} & (50, 20), 200 & :68-69 & 1 cm di Δs ≈ 2 cm di errore \\
\texttt{MPC\_S\_V}, \texttt{MPC\_W\_V} & (5000, 150), 2e-3 & :70-71 & quota rigida (nessuna molla verticale) \\
\texttt{DS\_MAX\_CAP} & 0,03 m & :72 & \textbf{vincolo sempre attivo}, θ ≤ 10,48° \\
\texttt{F\_MIN/MAX(\_THRUST)\_RATIO} & 0,3 / 3 / 6 × \(m_bg\) & :73-75 & contatto garantito; spinta \\
\texttt{VMC\_KP/KD} (appoggio) & (1500, 0) / (40, 15) & :81-82 & quota affidata a \(F_z\) \\
\texttt{KF\_Q\_ACC}, \texttt{KF\_R\_POS}, \texttt{KF\_R\_VEL} & 0,5; 1e-4; 1e-3 & MPC :89-91 & τ di fusione 45 / 280 ms \\
\bottomrule\noalign{}
\end{longtable}

\begin{longtable}[]{@{}>{\raggedright\arraybackslash}p{(\linewidth - 8\tabcolsep) * \real{0.2500}}>{\raggedright\arraybackslash}p{(\linewidth - 8\tabcolsep) * \real{0.2500}}>{\raggedright\arraybackslash}p{(\linewidth - 8\tabcolsep) * \real{0.2500}}>{\raggedright\arraybackslash}p{(\linewidth - 8\tabcolsep) * \real{0.2500}}@{}}
\toprule\noalign{}
\texttt{SagittalGains} (K\_s, K\_sd, ρ, k\_ξ, k\_i) & 60, 5, 0,4, 1,5, 0,5 & \texttt{rotino\_pid/zmp\_balance.py} & poli robusti sul VL-WIP discretizzato, §3.2 \\
\midrule\noalign{}
\endhead
\texttt{acc\_max} & 3 m/s² & \texttt{rotino\_pid/zmp\_balance.py} & offset ZMP ≤ 59 mm \\
\texttt{LEG\_KP\_STAND}, \texttt{LEG\_KD\_STAND} & (1500, 5000), (40, 80) & \texttt{rotino\_pid/controller.py} & nessun ciclo limite a 500 Hz, §3.4 \\
\texttt{K\_PSI}, \texttt{K\_OMEGA}, \texttt{K\_LAT} & 1,8 N·m/rad; 0,45 N·m·s/rad; 2,0 & \texttt{rotino\_pid/controller.py} & poli di imbardata −26,7, −4,7 \\
\bottomrule\noalign{}
\end{longtable}

\subsection{Appendice C. Riproducibilità}\label{appendice-c-riproducibilità}

Gli script \texttt{docs/verifiche/} citati nella prima versione di questo documento non sono nel repository: i numeri \textbf{[C]}, \textbf{[S-R]} e \textbf{[S-G]} di quella versione non sono quindi ripetibili così come sono. Per le parti aggiornate:

\begin{longtable}[]{@{}>{\raggedright\arraybackslash}p{(\linewidth - 4\tabcolsep) * \real{0.5000}}>{\raggedright\arraybackslash}p{(\linewidth - 4\tabcolsep) * \real{0.5000}}@{}}
\toprule\noalign{}
Comando & Cosa produce \\
\midrule\noalign{}
\endhead
\texttt{cd src/rotino\_pid \&\& python3 -m pytest -q test} & verifica della legge del PID sul VL-WIP lineare: robustezza, trapezio, spinta, anteprima LIPM \\
\texttt{ros2 run rotino\_benchmark suite} & tutti gli scenari con PID e MPC in Gazebo, poi \texttt{benchmark\_runs/suite\_<data>/riepilogo.md} \\
\texttt{ros2 run rotino\_benchmark campaign -- spinta} & un solo scenario, con \texttt{confronto.md} e grafici \\
\texttt{ros2 run rotino\_benchmark suite -- --list} & l'elenco degli scenari \\
\bottomrule\noalign{}
\end{longtable}

Nota: su questa macchina \texttt{scipy} è incompatibile con NumPy 2.2; gli script non lo usano.

\end{document}
